%======================================================================
\section{Part I --- Tensor-network notebook (\code{msh\_tjunctions\_tn.ipynb})}
%======================================================================

Every reported quantity in this notebook comes from MPO/MPS methods: DMRG for the subgap states and the Majoranas, TDVP for the
gates and the fusion readout. Dense matrices appear only in section 4, to check that the MPO equals the lattice Hamiltonian.
Packages: \code{ITensors}, \code{ITensorMPS}, \code{TensorBinding} (alias \code{TB}), \code{Plots}, \code{Serialization}.

\subsection{Setup (cell 1)}

Global constants: \code{ROOT = pwd()} (all files are read and written in the notebook's folder); \code{RUN\_TDVP = false} (set to
\code{true} to rerun the TDVP of section 6 even when the cache exists); the Pauli matrices \code{σ0, σx, σy, σz}.

\begin{function}{quiet(f)}
\begin{Verbatim}
quiet(f) = with_logger(f, ConsoleLogger(stderr, Logging.Error))
\end{Verbatim}
Runs the closure \code{f} with a logger that shows errors only. KrylovKit warns at every non-converged local solve in DMRG on $K$ and in
\code{linsolve}; the notebook comment calls these warnings harmless. Used around \code{zero\_mode\_subspace} and \code{refine\_subspace}.
\end{function}

%----------------------------------------------------------------------
\subsection{Section 1: Hamiltonian, MPO and dense reference}

\begin{function}{diag\_mpo\_exact(v, pos; cutoff=1e-14)}
\begin{Verbatim}
function diag_mpo_exact(v::AbstractVector, pos; cutoff=1e-14)
    A = permutedims(reshape(ComplexF64.(v), ntuple(_ -> 2, L)), L:-1:1)
    ψ = MPS(A, pos; cutoff=cutoff)
    ...  δ[s => j, s' => j, s'' => j] = 1
    W[k] = replaceprime(ψ[k] * δ, 2 => 0; tags=tags(s))
\end{Verbatim}
Builds the exact MPO of the diagonal operator $\sum_n v_n\ket n\bra n$ on the $L$ position qubits (pos$_1$ = most significant bit).
\emph{Method:} the length-$2^L$ vector is reshaped into an $L$-index tensor (the \code{permutedims} puts the most significant bit first),
compressed into an MPS by TT-SVD with \code{cutoff = 1e-14}, and each MPS tensor is contracted with a three-leg Kronecker delta
$\delta_{s s' s''}$ so that the physical index becomes a diagonal operator index pair. The result is exact up to the cutoff.
\emph{Used in:} the site-dependent exchange field (\code{msh\_tjunction\_array}, \code{array\_exchange\_mpo}) and the position operator
(\code{array\_position\_mpo}).
\end{function}

\begin{function}{dense(mpo, sites)}
Contracts the whole MPO (\code{prod}) and reshapes it into a $D\times D$ matrix, big-endian over \code{sites} (first site most
significant). Only used for the checks of section 4 and the \code{tb\_reference/} export.
\end{function}

\begin{function}{hop\_matrix(t, α, d)}
\begin{Verbatim}
hop_matrix(t, α, d) = -t * σ0 + im * α * (d[1] * σy - d[2] * σx)
\end{Verbatim}
The $2\times2$ spin matrix of the hopping $c^\dagger_{\mathbf r+\hat d}\,h\,c_{\mathbf r}$ along the unit bond $\hat d = (d_x, d_y)$:
$h = -t + i\alpha(d_x\sigma_y - d_y\sigma_x)$. The reverse bond has $h(-\hat d) = h(\hat d)^\dagger$.
\end{function}

\begin{function}{struct TArray; Ncell, center\_x, array\_pos, array\_spin, ARM\_ANG, ARM\_ORIG}
\begin{Verbatim}
Base.@kwdef struct TArray
    L::Int = 4          # 2^L = 16 sites per arm
    ncell::Int = 3
end
Ncell(a::TArray) = 2^a.L
center_x(a::TArray, c) = (c - (a.ncell + 1) / 2) * 2Ncell(a)
array_pos(a, c, s, n) = (center_x(a, c), 0.0) .+ ARM_ORIG[s] .+ n .* (cos(ARM_ANG[s]), sin(ARM_ANG[s]))
\end{Verbatim}
\code{TArray} holds the geometry: $L$ position qubits ($2^L$ sites per arm) and the number of cells. \code{Ncell} returns $2^L$ (sites per
arm, despite the name). \code{center\_x} gives the junction coordinate $x_c = (c - (n_{\rm cell}+1)/2)\cdot 2^{L+1}$, i.e.\ $-32, 0, 32$ for
$L = 4$. \code{array\_pos} returns $\mathbf r(c,s,n)$ with \code{ARM\_ANG = (0, π, π/2)} and \code{ARM\_ORIG = ((0,0), (-1,0), (0,1))}.
\code{array\_spin(a, f)} returns a closure $(c,s,n)\mapsto\mathbf S$ with $\theta = (-1)^{x+y}(\pi/2)f(c,s,n)$ and
$\phi = \theta_{\rm arm} + \pi/2$, using the rounded integer coordinates for the sign.
\end{function}

\begin{function}{msh\_tjunction\_array(a, f; t, μ, α, Δ, J, cutoff=1e-13, sites=nothing)}
\begin{Verbatim}
function msh_tjunction_array(a::TArray, f; t=1.0, μ=-7.0, α=0.9, Δ=2.4, J=5.8, cutoff=1e-13, sites=nothing)
    ...
    return HB, [ns; ss; cl; pos; arm]
end
\end{Verbatim}
Builds the BdG Hamiltonian of the array with texture \code{f} as an MPO and returns it with the site list
\code{[nambu; spin; cell; pos₁…pos₄; arm]}. If \code{sites} is given, the existing indices are reused, so MPS built on one Hamiltonian
can be used with another (needed for $H(t)$).

\emph{Construction.} A helper \code{term(Mpos, σmat, Bcell, Barm)} takes an MPO on the position qubits and attaches an operator on the
arm (\code{TB.postpend\_op}), the cell and the spin (\code{TB.prepend\_op}). \code{P(d,i)} is the projector $\ket i\bra i$ and
\code{E(d,i,j)} $= \ket i\bra j$. The normal part $H_0$ is the sum of
\begin{enumerate}
\item the chemical potential $-\mu\cdot\mathbb 1$;
\item intra-arm hopping for each arm $s$: $K_+\otimes h(\hat e_s)\otimes\mathbb1_{\rm cell}\otimes P_s$ and
  $K_-\otimes h(\hat e_s)^\dagger\otimes\mathbb1\otimes P_s$, where $(K_+, K_-)$ = \code{TB.shift\_pair\_mpos(pos, 1; cyclic=false)} are the
  binary shift $\ket n\to\ket{n+1}$ and its adjoint;
\item junction bonds through the site-0 projector \code{P0 = TB.\_site\_projector\_mpo(L, pos, 0)}, with arm operators $\ket b\bra a$ for the
  four directed bonds $(a\to b, \hat d)$: $(2\to1, +\hat x)$, $(1\to2, -\hat x)$, $(3\to1, -\hat y)$, $(1\to3, +\hat y)$. These connect the
  junction site (arm 1, $n=0$) with its left neighbour (arm 2, $n=0$) and with the first site of the vertical leg (arm 3, $n=0$);
\item inter-cell bonds through the site-15 projector \code{Pend}, with cell operator $\ket{c+1}\bra c$ and arm operator $\ket2\bra1$
  (right end of cell $c$ to the far end of cell $c+1$'s left arm) and the Hermitian conjugate;
\item the exchange $J\,\mathbf S\cdot\boldsymbol\sigma$: for each cell, arm and component $k$ the vector $v_n = JS_k(c,s,n)$ becomes a
  diagonal MPO (\code{diag\_mpo\_exact}), with $\sigma_k$ on spin and projectors on cell and arm; zero components are skipped.
\end{enumerate}
The terms are summed with \code{cutoff = 1e-13}. The BdG MPO is then
\[
H_B = \tau_{pp}\otimes H_0 + \tau_{hh}\otimes(-H_0^*) + \tau_+\otimes(\Delta\,i\sigma_y) + \tau_-\otimes(\Delta\,i\sigma_y)^\dagger,
\]
using \code{TB.prepend\_nambu} with \code{:Pp}, \code{:Ph}, \code{:tp} ($\ket p\bra h$), \code{:tm} ($\ket h\bra p$); $-H_0^*$ is
\code{-1.0 * dag(H0)} and the pairing is \code{ComplexF64[0 1; -1 0]} $= i\sigma_y$. The resulting MPO has $\chi = 18$.
\emph{Used in:} section 4 (\code{HB}), as the $J=0$ base of \code{MovingArrayH}, and in \code{bench\_tdvp}.
\end{function}

\begin{function}{array\_reference(a, f; t, μ, α, Δ, J)}
Independent dense BdG matrix ($576\times576$) built only from site coordinates, in the same basis as \code{dense(HB, sites)}. The index of
site $(c,s,n)$ is \code{idx = (c-1)*3N + 3n + s} (cell $\otimes$ pos $\otimes$ arm, arm fastest); spin blocks are \code{[i, M+i]}. On-site blocks
$-\mu + J\mathbf S\cdot\boldsymbol\sigma$; for every ordered pair of sites exactly one lattice spacing apart (\code{hypot(d) ≈ 1}),
$H_0[\mathbf r+\hat d,\mathbf r] \mathrel{+}= h(\hat d)$. This brute-force double loop does not use the arm/junction/cell logic of the MPO,
which makes it an independent check. Pairing block $\Delta\,i\sigma_y\otimes\mathbb1_M$; result $\begin{psmallmatrix}H_0 & D\\ D^\dagger & -H_0^*\end{psmallmatrix}$.
\end{function}

\begin{function}{array\_exchange\_mpo(a, sites, f, cells; J=5.8, cutoff=1e-14)}
The exchange of the listed \code{cells} only, as a BdG MPO on \code{sites}: particle part $\tau_{pp}\otimes\sigma_k$, hole part
$\tau_{hh}\otimes(-\sigma_k^*)$, each times the diagonal MPO of $JS_k$ and cell/arm projectors. Returns \code{nothing} if all components
vanish. This is what makes $H(t)$ cheap: only the exchange of moving cells is rebuilt at each time (\code{MovingArrayH}).
\end{function}

\begin{function}{array\_position\_mpo(a, sites; β=0.05)}
Diagonal position operator $\hat X = \sum_{\mathbf r} (\mathbf r\cdot\hat u)\ket{\mathbf r}\bra{\mathbf r}\otimes\mathbb1_{\rm spin}\otimes\mathbb1_{\rm nambu}$
with $\hat u = (\cos\beta, \sin\beta)$, $\beta = 0.05$. Built per cell and arm from \code{diag\_mpo\_exact}. Used by \code{localize\_majoranas};
the eigenvalues it produces are the ``centres $\mathbf r\cdot\hat u$'' printed in section 5.
\end{function}

\begin{function}{site\_list(a)}
All sites as $(c, s, n)$, ordered cell, then $n$, then arm. This is the row order of the LDOS CSV and of \code{coords.txt}.
\end{function}

%----------------------------------------------------------------------
\subsection{Section 2: protocols for \texorpdfstring{$\sqrt Z$, $\sqrt X$}{√Z, √X} and the partial fusion}

All protocols are textures $f(c,s,n;t)$ built from
\[
\mathrm{ss}(\tau) = \begin{cases}0 & \tau\le0\\ \sin^2(\pi\tau/2) & 0<\tau<1\\ 1 & \tau\ge1\end{cases}
\qquad(\code{smoothstep}),
\]
with $\tau = t/T_{\rm rot}$ and a site-dependent delay \code{step\_T}$\cdot d$: each site turns from topological to trivial (or back) over one
\code{rot\_T}, starting at a delay that depends on its position, so the topological/trivial domain wall moves along the leg. The
ramps were written for legs of $L_1 = L_2 = 14$ sites (a 1D index $i$ along the horizontal segment and $k$ along the vertical leg) and are
mapped onto the 16-site legs of the array.

\begin{function}{smoothstep(τ), f0}
\code{smoothstep} as above. \code{f0(c,s,n)} $=1$ if $c=2$ or $s=3$, else $0$: the initial setup with four MZMs.
\end{function}

\begin{function}{struct Exchange12; duration(p::Exchange12)}
\begin{Verbatim}
Base.@kwdef struct Exchange12
    L1::Int = 14; L2::Int = 14; rot_T::Float64 = 200.0; step_T::Float64 = 0.1
end
duration(p::Exchange12) = (p.step_T * (4p.L1 + 2p.L2 + 8) + 1) * p.rot_T
\end{Verbatim}
Parameters of \code{exchange\_12} (\code{rot\_Tz = 200}, \code{step\_Tz = 0.1}). Duration of one exchange
$T_{12} = (0.1\cdot92 + 1)\cdot200 = 2040$; the Z gate is two exchanges, $T_{\rm gate} = 4080$.
\end{function}

\begin{function}{f\_horizontal(p, i, t), f\_vertical(p, k, t)}
\code{f\_horizontal} is the rotation fraction of horizontal site $i$ of the original $L_1$-geometry: $i = 0\dots L_1$ is the left half,
$i = L_1+1$ the junction (always $0$, topological), $i = L_1+2\dots2L_1+2$ the right half; elsewhere 1. Each half is a difference of two
delayed smoothsteps, i.e.\ the site becomes trivial at one delay and topological again at a later one. \code{f\_vertical} is
$1 - \mathrm{ss}(\tau - s(k+L_1+2)) + \mathrm{ss}(\tau - s(L_2 - k + 3L_1 + L_2 + 5))$ for vertical site $k$: the leg becomes topological
from the junction outward and trivial again from the far end inward. Together they move $\gamma_1$ and $\gamma_2$ around the vertical leg.
\end{function}

\begin{function}{f\_braid(p, t), f\_array\_braid(a, p, f0, t)}
\begin{Verbatim}
f_braid(p::Exchange12, t) = (s, n) -> s == 1 ? f_horizontal(p, p.L1 + 1 + n, t) :
                                      s == 2 ? f_horizontal(p, max(p.L1 - n, 0), t) :
                                               f_vertical(p, min(n, p.L2), t)
\end{Verbatim}
The map onto the 16-site legs of cell 1: arm 1 $n\to i = L_1+1+n$ (the junction is $n=0$, $i = 15$; $n = 15\to i = 30 = 2L_1+2$);
arm 2 $n\to i = \max(L_1 - n, 0)$ (the two outermost sites $n = 14, 15$ share $i=0$); arm 3 $n\to k = \min(n, L_2)$ (sites $n = 14, 15$
share $k = 14$). \code{f\_array\_braid} applies this to cell 1 and \code{f0} elsewhere.
\end{function}

\begin{function}{struct Exchange23; duration; f\_h23, f\_v23, f\_cell2\_23, f\_array\_23}
\begin{Verbatim}
Base.@kwdef struct Exchange23
    L1::Int = 14; L2::Int = 14; rot_T::Float64 = 100.0; step_T::Float64 = 0.1
end
duration(p::Exchange23) = (p.step_T * (p.L1 + 3p.L2 + 2p.L1 + 3) + 1) * p.rot_T
\end{Verbatim}
\code{exchange\_23} acts on cell 2 (junction 2, trivial in the initial setup). Duration $(0.1\cdot87+1)\,T_{\rm rot}$: 970 for
\code{rot\_T = 100} (used in the Hamiltonian check) and 3880 for \code{rot\_T = 400} (used for the X gate, $T_{\rm gate} = 7760$).
\code{f\_h23(p, i, t)} is defined on junction 2's segment $i = 2L_1+3\dots4L_1+4$ ($31\dots60$) and raises an error elsewhere; each half
turns topological at one delay and trivial again at a later one, so the segment of cell 2 temporarily connects the two topological
segments and $\gamma_2,\gamma_3$ are exchanged through it. \code{f\_v23} has the same form as \code{f\_vertical}.
\code{f\_cell2\_23} maps the 16-site legs: right arm $i = 3L_1+3+n$ ($45\dots60$; junction $i = 45$); left arm
$i = (3L_1+2) - (L_1-1)\,n/(L_1+1)$, a linear map of the 16 sites onto $i = 44\dots31$ (non-integer $i$ is allowed since \code{smoothstep}
takes real arguments); vertical $k = \min(n, L_2)$. \code{f\_array\_23} applies it to cell 2 and \code{f0} elsewhere.
\end{function}

\begin{function}{struct Fusion; duration(q::Fusion); f\_fuse(q, f0, t)}
\begin{Verbatim}
Base.@kwdef struct Fusion
    L1::Int = 14; Lp::Int = 1; rot_T::Float64 = 100.0; step_T::Float64 = 0.2
end
duration(q::Fusion) = ((q.L1 - q.Lp + 1) * q.step_T + 1) * q.rot_T
\end{Verbatim}
Partial fusion of pair A on cell 1 (\code{rot\_T0 = 100}, \code{step\_T0 = 0.2}, \code{Lp = 1}), duration $T_f = 380$. In
\code{f\_fuse}: cell 1's vertical leg stays trivial; the sites $n < L_p$ of arm 2 and $n\le L_p$ of arm 1 stay topological (sites
$x_c-1, x_c, x_c+1$: $2L_p+1 = 3$ sites); every other site of cell 1 ramps to trivial as $\mathrm{ss}(\tau - 0.2(15-n))$, the outer end
$n = 15$ first. The leg length is hard-coded as \code{N = 16}. Both halves of the segment thus turn trivial from the outer ends inward.
\end{function}

%----------------------------------------------------------------------
\subsection{Section 3a: time evolution and readout}

\begin{function}{struct MovingArrayH; MovingArrayH(a, ft, f0, moving; J, sites); (B::MovingArrayH)(t)}
\begin{Verbatim}
struct MovingArrayH
    a::TArray; ft::Function; moving::Vector{Int}; J::Float64
    H_fixed::MPO; sites::Vector{<:Index}
end
(B::MovingArrayH)(t) = +(B.H_fixed, array_exchange_mpo(B.a, B.sites, B.ft(t), B.moving; J=B.J); cutoff=1e-14)
\end{Verbatim}
$H(t) = H_{\rm fixed} + H_{\rm ex}^{\rm moving}(t)$. The constructor builds $H_{\rm fixed}$ once: the full MPO at $J = 0$ (chemical
potential, hopping, Rashba, pairing) plus the exchange of the cells not in \code{moving}, computed with \code{f0}. Calling \code{B(t)}
adds the exchange of the moving cells for texture \code{ft(t)}. This avoids rebuilding the whole MPO at each time step.
\end{function}

\begin{function}{τz\_on(ψ), ΔQ\_mps(ψ)}
\code{τz\_on} applies $\tau_z = \mathrm{diag}(1,-1)$ to the Nambu site (the first MPS tensor) and leaves the rest unchanged.
\code{ΔQ\_mps(ψ)} $= \mathrm{Re}\bra\psi\tau_z\ket\psi = \sum(|u|^2 - |v|^2)$, the contribution of one mode to $\Delta Q$.
\end{function}

\begin{function}{site\_uv(ψ, sites, slist), site\_dn(ψ, sites, slist)}
\begin{Verbatim}
"Spin-summed particle |u|² and hole |v|² weights of an MPS at the sites `slist = [(c, s, n), …]` (no dense vector)."
...
        T = head * (ψ[3] * onehot(cl => c))
        for k in 1:L
            T *= ψ[3+k] * onehot(pos[k] => ((n >> (L - k)) & 1) + 1)      # pos[1] = most significant bit
        end
        A = Array(T * (ψ[end] * onehot(arm => s)), ns, ss)
\end{Verbatim}
Local weights of an MPS at selected sites, without forming the 576-component vector. For each site $(c,s,n)$ the cell, position bits
and arm indices are fixed with one-hot vectors; the remaining tensor \code{head} $=\psi[1]\psi[2]$ (Nambu $\times$ spin) gives the
$2\times2$ amplitude array $A_{\tau\sigma}$. Then $|u|^2 = \sum_\sigma|A_{1\sigma}|^2$, $|v|^2 = \sum_\sigma|A_{2\sigma}|^2$.
\code{site\_dn} returns $|u|^2 - |v|^2$, the contribution of one mode to $\delta n$. \emph{Used in:} the LDOS weights (section 5,
$|u|^2+|v|^2$) and $\delta n(x,t)$ in \code{evolve\_mode\_tdvp}.
\end{function}

\begin{function}{mode\_leak(z)}
\code{mode\_leak(z) = [1 - sum(abs2, z[:, k]) for k in 1:size(z, 2)]} for the $4\times2$ overlap matrix
$z_{ik} = \langle\gamma_i(0)|\psi_k\rangle$: the leakage $\ell_k$ of each evolved mode out of the initial zero-mode subspace. This presumes
the $\gamma_i$ orthonormal, which \code{clean\_majoranas} enforces.
\end{function}

\begin{function}{evolve\_mode\_tdvp(B, ψ0, tsnap; γ0, slist, dt=0.1, log, maxdim=64, cutoff=1e-20, tol=1e-12, krylovdim=30)}
\begin{Verbatim}
function evolve_mode_tdvp(B, ψ0, tsnap; γ0, slist, dt=0.1, log=(args...) -> nothing,
                          maxdim=64, cutoff=1e-20, tol=1e-12, krylovdim=30)
    ...
        n = round(Int, (t1 - t0) / dt); h = (t1 - t0) / n
        ψ = last(TB.evolve_with_tdvp_timedep(t -> B(t0 + t), ψ, n, h; maxdim, cutoff, tol, krylovdim, reverse_step=true))
        record!(i, ψ)
\end{Verbatim}
Evolves one mode $\psi(t)$ with $i\partial_t\psi = H(t)\psi$ from snapshot to snapshot. Between consecutive snapshot times $t_0, t_1$ it takes
$n = \mathrm{round}((t_1-t_0)/dt)$ steps of length $h = (t_1-t_0)/n$ with \code{TB.evolve\_with\_tdvp\_timedep}: 2-site TDVP in which, on each
interval $[t, t+h]$, $H$ is frozen at the midpoint $t+h/2$ (a second-order midpoint rule in time), states are normalized each step,
and \code{reverse\_step = true} (the backward local evolution of the TDVP sweep). Settings: \code{maxdim = 64}, \code{cutoff = 1e-20} (so
the bond dimension is set by the state, not by truncation; the runs reach $\chi = 24$), Krylov exponentiation with \code{tol = 1e-12},
\code{krylovdim = 30}.

At every snapshot \code{record!} stores $z_{i} = \langle\gamma_i(0)|\psi\rangle$ for the four initial Majoranas, \code{ΔQ\_mps(ψ)},
\code{site\_dn} on \code{slist} and \code{maxlinkdim(ψ)}; then \code{log(i, nt, t, χ, ΔQ)} is called.
\emph{Returns} \code{(; ts, z, dQ, dn, χ, ψT)}; \code{ψT} is the final MPS.
\end{function}

%----------------------------------------------------------------------
\subsection{Section 3b: zero modes from DMRG, refinement and localization}

The zero modes cannot be found by plain DMRG on $H$ (they are in the middle of the spectrum). The notebook minimises
$K = H^2 + \eta^2$, whose lowest states are the $|E|$-smallest states of $H$.

\begin{function}{norm\_Hψ(H, ψ)}
$\|H\psi\|$ by exact MPO--MPS application (\code{alg="naive"}, \code{truncate=false}); the residual used to judge the zero modes.
\end{function}

\begin{function}{zero\_mode\_subspace(HB, sites; nzero=2, η=1e-2, seed=1234, nsweeps=20, maxdim=[10,20,50,100,200], cutoff=1e-14, cutoff\_K=1e-14, noise=[1e-6,1e-7,1e-8,0], weight=10, eigsolve\_krylovdim=20, eigsolve\_maxiter=3)}
\begin{Verbatim}
    K = TB.build_K(HB, sites, 0.0, η; cutoff_K)        # TB's default cutoff_K = 1e-8 would drop η²
    for k in 1:nzero+1
        ψ0 = random_mps(sites; linkdims=10)
        _, ψ = isempty(ψs) ? dmrg(K, ψ0; ...) : dmrg(K, ψs, ψ0; weight, ...)
\end{Verbatim}
Finds \code{nzero} states of smallest $|E|$ plus one extra state. \emph{Method:} \code{TB.build\_K(H, sites, ω=0, η)} builds
$K = (H-\omega)^2 + \eta^2$ as an MPO (product of MPOs, truncated at \code{cutoff\_K}). The cutoff is set to $10^{-14}$ because, per the
comment, TensorBinding's default $10^{-8}$ would drop $\eta^2 = 10^{-4}$. DMRG then finds the ground state of $K$ and, one at a time, the
next states with a penalty \code{weight}$\,\sum_j\ket{\psi_j}\bra{\psi_j}$ on the states already found (ITensorMPS excited-state DMRG). The
bond dimension grows along \code{maxdim}; \code{noise} helps the early sweeps. \code{eigsolve\_krylovdim = 20} is needed: the markdown
states that the default stalls at $\|H\psi\|\approx0.2$.

The last state is removed and used only as a gap check: its $\|H\psi\|$ is the edge of the window (0.268 in section 5).
\emph{Returns} \code{(ψs, [‖Hψ‖ for each], ‖Hψ\_gap‖)}.
\end{function}

\begin{function}{refine\_subspace(HB, ψs; σ=1e-3, nsweeps=6, cutoff=1e-20, maxdim=200, krylovdim=30)}
\begin{Verbatim}
refine_subspace(HB, ψs; σ=1e-3, ...) =
    [(x = ITensorMPS.linsolve(HB, ψ, ψ, -σ, 1.0; nsweeps, cutoff, maxdim,
                              updater_kwargs=(; krylovdim, verbosity=0)); x / norm(x)) for ψ in ψs]
\end{Verbatim}
One inverse-iteration step for each state: solves $(H - \sigma)\,x = \psi$ by DMRG-like sweeps (\code{linsolve} with $a_0 = -\sigma$,
$a_1 = 1$) and normalizes. Components with $|E|\ll$ gap are amplified by $\sim|E_{\rm bulk}/\sigma|$ relative to bulk admixtures, which the
markdown identifies as what $\Delta Q$ is sensitive to. The comment records the trade-off for \code{krylovdim}: the local GMRES solves rarely
converge at 30; at 60 they do ($\|H\psi\|$ $2\times10^{-5}\to3\times10^{-8}$) but the step is about 2.5 times slower.
\end{function}

\begin{function}{mps\_comb(ψs, c; cutoff=1e-14)}
Normalized MPS linear combination $\sum_j c_j\psi_j/\|\cdot\|$, summed pairwise with \code{cutoff}.
\end{function}

\begin{function}{ritz(HB, ψs)}
Rayleigh--Ritz of $H$ in $\mathrm{span}(\psi_s)$: builds the overlap matrix $S_{ij} = \langle\psi_i|\psi_j\rangle$ and
$H_{ij} = \langle\psi_i|H|\psi_j\rangle$, solves the generalized Hermitian problem $Hc = ESc$, and returns the energies (ascending) and the
Ritz vectors as MPS. This separates the $\pm E$ members, which DMRG on $K = H^2+\eta^2$ cannot distinguish.
\end{function}

\begin{function}{localize\_majoranas(HB, sites, ψs, X)}
\begin{Verbatim}
    F  = eigen(Hermitian((Xm + Xm') / 2), Hermitian((S + S') / 2))
    τx = ITensor([0.0 1.0; 1.0 0.0], ns', ns)
    C(ψ) = (φ = dag(ψ); φ[1] = noprime(τx * φ[1]); φ)
    fix(γ) = exp(im * angle(inner(γ, C(γ))) / 2) * γ
\end{Verbatim}
Localizes the zero-mode subspace with the position MPO $X$ (\code{array\_position\_mpo}): diagonalizes $X$ in the subspace,
$X_{ij}c = r\,S_{ij}c$, so each eigenvector is the combination with a sharp position; the eigenvalues are the Majorana centres. It defines
the particle--hole conjugation $C\psi = \tau_x\psi^*$ (complex conjugate the MPS, $\tau_x$ on the Nambu site), and fixes the global phase of
each $\gamma$ so that $\langle\gamma|C\gamma\rangle$ is real and positive. \emph{Returns} \code{(; γs, centers, Hγ, C)}.
\end{function}

\begin{function}{clean\_majoranas(γs, C)}
\begin{Verbatim}
"""
Make each γ exactly self-conjugate, γ ← (γ + Cγ)/‖·‖, then orthonormalize with the real Löwdin coefficients S^(-1/2)
(⟨γᵢ|γⱼ⟩ is real for self-conjugate γ, so real combinations keep Cγ = γ). The leakage of section 6 assumes both.
"""
\end{Verbatim}
Projects each $\gamma$ onto the self-conjugate subspace, then Löwdin-orthonormalizes with $W = S^{-1/2}$ ($S$ real symmetric), so that the
$\gamma_i$ are exactly Majorana and orthonormal. Result in section 5: $\max\|C\gamma-\gamma\| = 1.3\times10^{-15}$,
$\max|\langle\gamma_i|\gamma_j\rangle - \delta_{ij}| = 1.8\times10^{-9}$.
\end{function}

\begin{function}{pair\_modes(HB, γs)}
Forms $\psi_A = (\gamma_1 + i\gamma_2)/\sqrt2$ and $\psi_B = (\gamma_3+i\gamma_4)/\sqrt2$. If $\langle\psi|H|\psi\rangle < 0$ it flips the sign of
$\gamma_2$ (or $\gamma_4$); for self-conjugate $\gamma$ this replaces $\psi$ by its particle--hole partner $C\psi = (\gamma_1 - i\gamma_2)/\sqrt2$,
which has the opposite energy. Both returned modes therefore have
$\langle\psi|H|\psi\rangle\ge0$. \emph{Returns} the (possibly sign-flipped) \code{γs} and \code{[ψ\_A, ψ\_B]}. The \code{γs} are sorted by
their position (from \code{localize\_majoranas}), so pairs (1,2) and (3,4) are the two ends of cell 1's and cell 3's segments.
\end{function}

%----------------------------------------------------------------------
\subsection{Section 4: check of the MPO and export of \code{tb\_reference/}}

The cell builds \code{a = TArray(L=4, ncell=3)}, \code{HB, sites = msh\_tjunction\_array(a, f0)} and compares with
\code{array\_reference}; it also checks particle--hole symmetry with $C_d = \sigma_x\otimes\mathbb1$. Output:

\begin{Verbatim}
D = 576, χ(H) = 18, ‖H_MPO − H_dense‖ = 3.5e-12, PH ‖C H* C + H‖ = 0.0e+00
exchange_12  t = 1020: ‖H_MPO(t) − H_dense(t)‖ = 4.3e-10, χ = 22
exchange_23  t =  485: ‖H_MPO(t) − H_dense(t)‖ = 2.0e-09, χ = 21
fusion       t =  300: ‖H_MPO(t) − H_dense(t)‖ = 1.9e-11, χ = 21
\end{Verbatim}
The protocol points are mid-exchange ($t = 1020 = T_{12}/2$; $t = 485 = T_{23}/2$ at \code{rot\_T = 100}) and late in the fusion ($t = 300$),
each via \code{MovingArrayH}.

\begin{function}{export\_tb\_reference(a, H\_at; out=joinpath(ROOT, "tb\_reference"))}
Writes \code{coords.txt} (integer $x\ y$ of each site, order cell, $n$, arm) and the dense matrices \code{initial.bin},
\code{z\_t1020.bin}, \code{x100\_t485.bin}, \code{fuse\_t300.bin} (raw \code{ComplexF64}, Julia column-major). Read by section 4 of
\code{tjunctions\_python\_ed.ipynb}.
\end{function}

%----------------------------------------------------------------------
\subsection{Section 5: subgap LDOS and MZMs from DMRG}

\begin{function}{gauss(x); constants ηg, ωs}
$g_\eta(x) = e^{-x^2/2\eta_g^2}/(\sqrt{2\pi}\eta_g)$ with $\eta_g = 0.01$; frequency grid \code{ωs} = 401 points on $[-0.2, 0.2]$.
\end{function}

\paragraph{Steps of the section (top-level code).}
\begin{enumerate}
\item \code{zero\_mode\_subspace(HB, sites; nzero=6)}: the six smallest-$|E|$ states (four MZMs and the $\pm0.127$ bound state of the trivial
  junction 2) plus the gap-check state; 41~s.
\item \code{ritz}: energies $[-0.127, -1.12\times10^{-9}, -1.08\times10^{-9}, 1.03\times10^{-9}, 1.13\times10^{-9}, 0.127]$, window edge 0.268.
\item LDOS $A(\mathbf r,\omega) = \sum_k w_k(\mathbf r)\,g_\eta(\omega - E_k)$ with $w_k = |u|^2+|v|^2$ from \code{site\_uv}: total
  (particle + hole, spin-summed). The window $|\omega|\le0.2$ lies inside the edge 0.268, so these six states are all that contribute.
  Check: $\sum_{\mathbf r}\int A\,d\omega = 6.0000$; max $A = 17.34$. Written to \code{msh\_array\_ldos\_tn.csv} (columns \code{x,y,omega,A}).
\item MZMs: the four Ritz states of smallest $|E|$ $\to$ \code{refine\_subspace} ($\sigma = 10^{-3}$) $\to$ \code{localize\_majoranas} with
  \code{array\_position\_mpo} $\to$ \code{clean\_majoranas} $\to$ \code{pair\_modes}, giving \code{γ0} and \code{ψ0}; 279~s. Centres
  $-46.5, -17.5, 16.5, 45.5$; $\langle\psi|H|\psi\rangle = 1.09\times10^{-9}, 1.16\times10^{-9}$; $\Delta Q$ of the two modes
  $2.4\times10^{-6}$ and $-2.5\times10^{-7}$.
\end{enumerate}

\begin{function}{arrows!(p, x, y, u, v; ℓ=0.85, head=0.38, hw=0.42, color, lw, label)}
Plot helper: draws arrows of length $\ell$ centred on $(x,y)$ along $(u,v)$ in data coordinates (a line for the shaft, a filled triangle for the
head, NaN-separated so one series draws all arrows). Used for the in-plane moments in Fig.~\ref{fig:ldos}.
\end{function}

%----------------------------------------------------------------------
\subsection{Section 6: gates followed by partial fusion with TDVP}

Three runs, each one evolution of $\psi_A$ and $\psi_B$ from \code{ψ0}:
\begin{center}
\begin{tabular}{llllll}
\toprule
Run & Gate & Moving cells & $T_1$ & $T_{\rm gate}$ & $T_{\rm end}$\\
\midrule
Z & \code{exchange\_12} twice, then fusion & [1] & 2040 & 4080 & 4460\\
X & \code{exchange\_23} (\code{rot\_T = 400}) twice, then fusion & [1, 2] & 3880 & 7760 & 8140\\
F & fusion only & [1] & 0 & 0 & 380\\
\bottomrule
\end{tabular}
\end{center}
Snapshots every 5 time units during the gate and every 1 during the fusion. For X both cells move: cell 2 during the gate and cell 1
during the fusion.

\paragraph{Horizontal chain.} \code{hsites} lists the 96 sites of $y=0$ from left to right (per cell: arm 2 $n = 15\dots0$, then arm 1
$n = 0\dots15$); \code{xh} their $x$ coordinates. $\delta n$ is recorded on these sites.

\begin{function}{protocol(kind)}
For \code{"Z"}, \code{"X"} or \code{"F"} returns \code{(; B, T1, Tg, tsnap)}: a \code{MovingArrayH} with texture
\code{t < Tg ? fg(t < T1 ? t : t - T1) : f\_fuse(q, f0, t - Tg)} (the second exchange repeats the first), the duration of one exchange,
the gate duration $T_g = 2T_1$, and the snapshot times \code{0:5:Tg} followed by \code{Tg+1:1:Tg+Tf}.
\end{function}

\begin{function}{tdvp\_log(key, msg; append=true), stamp(), hms(s); LOG\_FILE, STATUS\_FILE}
Thread-safe progress reporting (a \code{ReentrantLock}). \code{tdvp\_log} appends a time-stamped line to \code{tdvp\_progress.log} (if
\code{append}), stores the message as the latest status of evolution \code{key} (e.g.\ \code{"ZA"}), and rewrites \code{tdvp\_status.txt}
with one line per evolution. \code{stamp()} gives the date and time, \code{hms(s)} formats seconds as h:mm:ss.
\end{function}

\begin{function}{run\_mode(kind, k, P)}
Runs \code{evolve\_mode\_tdvp(P.B, ψ0[k], P.tsnap; γ0, slist=hsites, log=logf)} for mode $k$ (A or B) of run \code{kind}. The inner
\code{logf} writes to the status file at every snapshot and to the log at every 1\% of $T_{\rm end}$, with $t$, $\chi$, the mode's $\Delta Q$,
elapsed time and an estimate of the remaining time (elapsed$\cdot(1-f)/f$). Errors are logged as \code{FAILED} and rethrown. Returns the
result merged with the wall time \code{secs}.
\end{function}

\paragraph{Final Hamiltonian.} \code{HT = protocol("F").B(Tf)} is the Hamiltonian at the end of the fusion (the same for every run).
\code{zero\_mode\_subspace(HT, sites; nzero=4)}, \code{refine\_subspace} and \code{ritz} give its four lowest states
\code{φF} with energies $[-0.0439, -1.16\times10^{-9}, 1.16\times10^{-9}, 0.0439]$ (253~s). Pair A is split to $\pm E_A$, pair B stays at zero.

\begin{function}{fusion\_leak(ψ)}
\begin{Verbatim}
fusion_leak(ψ) = (c = [inner(φ, ψ) for φ in φF]; 1 - real(c' * (SF \ c)))
\end{Verbatim}
Leakage of an evolved mode out of $\mathrm{span}(\phi_F)$: $1 - c^\dagger S_F^{-1}c$ with $c_i = \langle\phi_i|\psi\rangle$ and the Gram matrix
$S_F$, i.e.\ the norm of the projection onto a possibly non-orthogonal basis. The adiabatic $\Delta Q$ of pair A is
\code{ΔQ\_mps(φF[iA])} for the $+E_A$ state \code{iA = argmax(EF)}: $-0.476854$.
\end{function}

\begin{function}{cache(kind); threaded execution}
\code{cache(kind)} is the path \code{msh\_tjunctions\_tdvp\_\{kind\}.jls}. Runs whose cache is missing (or all, if \code{RUN\_TDVP}) are
started as \code{@spawn(run\_mode(kind, k, P))} tasks, two per run, so up to six evolutions run in parallel, one per Julia thread, with
BLAS set to one thread. For each run the results of both modes are serialized: \code{ts}, \code{z}, \code{dQ}, \code{dn}, \code{χ},
\code{leak\_fus = fusion\_leak(ψT)}, \code{secs}, \code{threads}. Only numbers are cached, not the MPS. In the saved notebook the cache
existed (the kernel ran on 1 thread); the timings in the cache are from the six-thread run.
\end{function}

\begin{function}{analyze(kind)}
Loads the cache and returns a \code{Dict} with \code{:Tg}, \code{:ts}, \code{:dQ} (sum of both modes), \code{:D} ($\delta n$, sum of both
modes, snapshot $\times$ site), \code{:χ} (max of both), \code{:secs}, \code{:threads}; for gate runs \code{:leak\_one} and \code{:leak}
(\code{mode\_leak} at $t = T_1$ and $t = T_g$); and \code{:fusion} with the final $\Delta Q$, the per-mode $\Delta Q$, the fusion leakage and the
time trace $\Delta Q(t - T_g)$ during the fusion.
\end{function}

\begin{function}{write\_density(name, r), density\_figure(ts, D, Tg, gate; method)}
\code{write\_density} writes \code{msh\_gate\_density\_tdvp\_\{name\}.csv} (\code{t,x,dn}, $x$ fastest) and
\code{msh\_gate\_dQ\_tdvp\_\{name\}.csv} (\code{t,dQ}). \code{density\_figure} draws two panels: $\log_{10}|\delta n(x,t)|$ over the whole
protocol (clipped at $10^{-16}$, colour range $[-12, 0]$; dashed lines at $T_g/2$ and $T_g$, dotted at the junctions) and $\delta n(x,t)$ during
the fusion on a symmetric colour scale, with $\sum_{\rm chain}\delta n$ at the end in the title.
\end{function}

\paragraph{Results.} Z: leakage after one exchange $(9.72\times10^{-7}, 1.65\times10^{-11})$, after two $(1.99\times10^{-6},
2.07\times10^{-11})$; fusion $\Delta Q = -0.477717$ (A $-0.477717$, B $1.4\times10^{-7}$), leak $(3.69\times10^{-5}, 1.44\times10^{-11})$.
Fusion only: $\Delta Q = -0.477372$, leak $(3.78\times10^{-5}, 7.0\times10^{-12})$. X: leakage after two exchanges
$(3.42\times10^{-5}, 3.49\times10^{-5})$; fusion $\Delta Q = +0.477004$ (A $+0.476113$, B $+8.9\times10^{-4}$), leak $(8.08\times10^{-5},
3.33\times10^{-5})$. The sum over the horizontal chain at the end is $\mp0.444$ (figure titles); the rest of $\Delta Q$ lies on the vertical leg,
since $\Delta Q$ is summed over all sites while $\delta n$ is recorded on $y=0$ only.

\begin{figure}[H]
\centering
\includegraphics[width=\textwidth]{tn_density_Z.png}
\caption{TN, Z run (\code{density\_figure}). Left: $\log_{10}|\delta n(x,t)|$ on the horizontal chain over two \code{exchange\_12} and the
fusion (dashed: $T_1 = 2040$, $T_g = 4080$). The weight of $\psi_A$ moves around junction 1 ($x=-32$) while cell 3 ($x>16$) is static.
Right: $\delta n$ during the fusion; the charge appears at the three topological sites around junction 1.}
\end{figure}

%----------------------------------------------------------------------
\subsection{Section 7: \texorpdfstring{$\Delta Q$}{ΔQ} during the fusion and bond dimension}

Top-level plotting code (no new functions). Left panel: $\Delta Q(t-T_g)$ for Z, X and F; right: $\max\chi(\psi_A,\psi_B)$ against
$t/T_{\rm end}$, with the start of the fusion marked. The bond dimension saturates at 24 early in every run, well below \code{maxdim = 64}.

\begin{figure}[H]
\centering
\includegraphics[width=0.95\textwidth]{tn_dQ_chi.png}
\caption{TN, section 7. Left: $\Delta Q$ during the partial fusion; Z and no gate give $-0.477$, X gives $+0.477$. The charge develops for
$t - T_g\gtrsim 260$, when the last ramped sites of cell 1 (arm $n = 2$, delay $0.2\cdot13$, i.e.\ $t-T_g = 260$) turn trivial. Right: largest bond dimension of the two modes, saturating at $\chi=24$.}
\end{figure}

%----------------------------------------------------------------------
\subsection{Section 8: summary}

Top-level code writes 12 \code{key,value} rows to \code{msh\_tjunctions\_tn\_gates.csv}: \code{Z\_leak\_A/B}, \code{X\_leak\_A/B} (after two
exchanges), \code{Z/X/F\_fus\_dQ}, \code{Z/X/F\_fus\_leak\_A}, \code{fus\_dQad\_A} ($-0.476854$) and \code{fus\_E\_A} ($0.0438539$). The
values appear in Table~\ref{tab:compare}.

%----------------------------------------------------------------------
\subsection{Section 9: cost of the tensor-network route}

Measured wall times (from the cache for TDVP; each evolution on one thread, six at once):

\begin{center}
\begin{tabular}{lrrr}
\toprule
Step & Wall time & Steps & Per step\\
\midrule
DMRG, six subgap states & 0:00:41 & & \\
refinement + localization & 0:04:39 & & \\
DMRG, final Hamiltonian & 0:04:13 & & \\
TDVP Z, mode A / B & 1:17:48 / 1:06:57 & 44\,600 & 104.7 / 90.1 ms\\
TDVP X, mode A / B & 2:18:53 / 2:16:30 & 81\,400 & 102.4 / 100.6 ms\\
TDVP F, mode A / B & 0:08:06 / 0:06:32 & 3\,800 & 127.9 / 103.3 ms\\
\bottomrule
\end{tabular}
\end{center}

\begin{function}{bench\_tdvp(L; χ=24, nsteps=10)}
\begin{Verbatim}
"Time per TDVP step (dt = 0.1, maxdim = 64) for arms of 2^L sites, initial Hamiltonian, random MPS of bond dimension χ."
\end{Verbatim}
Builds the array with arms of $2^L$ sites, a \code{MovingArrayH} with the static texture \code{f0}, and a random complex MPS of bond dimension
24 (the value reached in section 6). One step is run first to exclude compilation, then \code{nsteps} are timed with the same TDVP settings
as section 6. \emph{Returns} \code{(; L, nsite, per\_step, mps\_size, dense\_size)}, where \code{mps\_size} is the number of complex entries of
the MPS and \code{dense\_size} $= 4\,n_{\rm site}$.

\begin{center}
\begin{tabular}{rrrrr}
\toprule
$L$ & sites & TDVP step & MPS entries & dense entries\\
\midrule
4 & 144 & 93.3 ms & 4845 & 576\\
5 & 288 & 115.4 ms & 5997 & 1152\\
6 & 576 & 247.3 ms & 7149 & 2304\\
\bottomrule
\end{tabular}
\end{center}
Doubling the number of sites adds one qubit to the MPS, so its size grows by a constant ($\approx1150$ entries) per doubling. All times and
sizes go to \code{msh\_tjunctions\_tn\_timing.csv}.
\end{function}
